\subsubsection{Subsonic Outflow Boundary Condition}%
\label{SEC:BC:Outflow:Subsonic}

This boundary condition assumes the velocity normal to the boundary face is
subsonic and leaving the computational domain.
%
The right state is computed as
%
\begin{alignat*}{1}
  \vec{u}_\subR & = \vec{u}_\subL \\
  T_\subR & = T_\subL \\
  p_\subR & = p_\subBC
\end{alignat*}
%
This boundary condition does have some built-in safe guards in case the
velocity at the boundary face becomes supersonic or reversed (\thatis inflow
instead of outflow).
%
However, these safe guards only provide minimal protection against
minor violations of the assumed subsonic outflow.
%
\begin{itemize}[itemsep=5pt,topsep=5pt,parsep=0pt]
  \item If you suspect your outflow boundary will likely contain large sections
    of both subsonic and supersonic outflow, you should consider using either
    the general outflow (\ref{SEC:BC:Outflow:General}) or characteristic
    (\ref{SEC:BC:Farfield:Characteristic}) boundary conditions.
  \item If you suspect there is a chance parts of your outflow boundary could
    reverse direction, at a minimum you should switch to using a characteristic
    (\ref{SEC:BC:Farfield:Characteristic}) boundary condition instead of a
    subsonic outflow boundary condition. However, the better option is to modify
    the grid so that the outflow boundary is in a region that maintains a flow
    directed out of the domain and is not susceptible to reversing that
    direction.
\end{itemize}

%###############################################################################

